% ---- pg-546 （书页 534） ----

我们首先搜索各对语言(聚类)之间的最小距离. 最小距离 1 出现在丹麦语与挪威语, 意大利语与法语, 及意大利语与西班牙语之间. 将语言依出现在表中的顺序编号, 我们有
\[
d_{32}=1;\qquad d_{86}=1;\qquad d_{87}=1
\]
由于 $d_{76}=2$, 只能合并 8 和 6, 或者合并 8 和 7. 但不能在 1 的水平上将 6,7,8 三者合并. 我们选择先合并 6 和 8, 然后更新距离阵, 并合并 2 和 3, 得到聚类(68)和(23). 随后的计算机计算产生\hyperref[fig:12.4]{图 12.4} 中的树图.

\begin{mvfigure}
\begin{tikzpicture}[x=0.92cm,y=0.42cm,line width=0.5pt]

% ---------- 叶结点的竖线 ----------
\draw (0,0) -- (0,2.2);
\draw (1,0) -- (1,0.87);
\draw (2,0) -- (2,0.87);
\draw (3,0) -- (3,0.87);
\draw (4,0) -- (4,0.87);
\draw (5,0) -- (5,2.16);
\draw (6,0) -- (6,5.05);
\draw (7,0) -- (7,5.05);
\draw (8,0) -- (8,5.05);
\draw (9,0) -- (9,8.64);
\draw (10,0) -- (10,8.64);

% ---------- 合并的横线 ----------
\draw (1,0.87) -- (2,0.87);
\draw (3,0.87) -- (4,0.87);
\draw (0,2.2) -- (1.5,2.2);
\draw (3.5,2.16) -- (5,2.16);
\draw (4.25,5.05) -- (6,5.05);
\draw (7,5.05) -- (8,5.05);
\draw (0.75,7.65) -- (5.125,7.65);
\draw (9,8.64) -- (10,8.64);
\draw (2.94,9.78) -- (9.5,9.78);

% ---------- 各聚类向上的竖线 ----------
\draw (1.5,0.87) -- (1.5,2.2);
\draw (3.5,0.87) -- (3.5,2.16);
\draw (0.75,2.2) -- (0.75,7.65);
\draw (4.25,2.16) -- (4.25,5.05);
\draw (5.125,5.05) -- (5.125,7.65);
\draw (7.5,5.05) -- (7.5,9.78);
\draw (2.94,7.65) -- (2.94,9.78);
\draw (9.5,8.64) -- (9.5,9.78);

% ---------- 坐标轴 ----------
\draw (-0.85,0) -- (-0.85,10.4);
\foreach \t in {0,2,4,6,8,10}{%
  \draw (-0.85,\t) -- (-0.62,\t);%
  \node[left,font=\small] at (-0.95,\t) {\t};%
}
\node[rotate=90,font=\small] at (-1.75,5) {距离};
\node[font=\small] at (5,-2.1) {语言};

% ---------- 叶结点的圆点与标记 ----------
\foreach \i/\l in {0/E,1/N,2/Da,3/Fr,4/I,5/Sp,6/P,7/Du,8/G,9/H,10/Fi}{%
  \fill (\i,0) circle (0.055);%
  \node[below,font=\small,inner sep=1.5pt] at (\i,0) {\l};%
}

\end{tikzpicture}
\phantomsection\label{fig:12.4}\mvcap{图 12.4\quad 11 种语言中前 10 个数词之间距离的单连接树图}
\end{mvfigure}

由这个树图我们看到, 挪威语和丹麦语, 以及法语和意大利语在最小距离(最大相似)水平上形成聚类, 当所允许的距离增大时, 英语加入挪威语—丹麦语组, 而西班牙语则并入法语—意大利语组. 注意匈牙利语和芬兰语之间的相似性比它们与其他语言聚类之间的相似性更大. 然而这两个聚类(语言)要到最后邻近点之间的距离有了很大增加之后才能合并. 最后, 所有的语言聚类在最近邻近点距离的最大值 9 这个水平上合并成一个聚类. \hfill$\blacksquare$

由于单连接方法用聚类之间的最短连接来合并聚类, 所以这种方法不能识别分离得很差的聚类[见\hyperref[fig:12.5]{图 12.5}(a)]. 另一方面, 单连接法是最能描绘非椭球形聚类的少数方法之一. 单连接法有一种挑选长链状聚类的倾向, 称为\mvbold{链接倾向}[见\hyperref[fig:12.5]{图 12.5}(b)]. 当项目位于链的两端因而颇不相似时, 链接法有可能产生误导.

\begin{mvfigure}
\begin{tikzpicture}[x=1.05cm,y=0.52cm,line width=0.5pt]

% ============ 点云宏: cx, cy, a, b, 倾角 ============
\def\cloud#1#2#3#4#5{%
  \pgfmathsetmacro\cth{cos(#5)}%
  \pgfmathsetmacro\sth{sin(#5)}%
  \foreach \f in {0.35,0.65,0.92}{%
    \foreach \a in {0,30,...,330}{%
      \pgfmathsetmacro\px{#1 + \f*(#3*cos(\a)*\cth - #4*sin(\a)*\sth)}%
      \pgfmathsetmacro\py{#2 + \f*(#3*cos(\a)*\sth + #4*sin(\a)*\cth)}%
      \fill (\px,\py) circle (0.04);%
    }%
  }%
}

% ================= (a) =================
\begin{scope}[shift={(0,0)}]
\draw (0,0) -- (3.4,0);
\draw (0,0) -- (0,3.4);
\node[rotate=90,font=\small] at (-0.45,1.7) {变量2};
\node[font=\small] at (1.7,-0.42) {变量1};
\cloud{1.15}{2.15}{0.85}{0.50}{35}
\cloud{2.15}{1.05}{0.80}{0.50}{35}
\node[font=\small,anchor=west] at (2.35,2.15) {椭圆形结构};
\node[font=\small] at (1.7,-1.05) {(a) 单连接在接近重叠时产生混淆};
\end{scope}

% ================= (b) =================
\begin{scope}[shift={(4.6,0)}]
\draw (0,0) -- (3.4,0);
\draw (0,0) -- (0,3.4);
\node[rotate=90,font=\small] at (-0.45,1.7) {变量2};
\node[font=\small] at (1.7,-0.42) {变量1};
% 链状连接
\coordinate (A) at (0.55,2.35);
\coordinate (B) at (1.05,1.75);
\coordinate (C) at (1.70,1.30);
\coordinate (D) at (2.30,1.45);
\coordinate (E) at (2.75,1.95);
\coordinate (F) at (0.85,2.75);
\coordinate (G) at (2.05,2.40);
\coordinate (H) at (2.60,0.95);
\draw (A) -- (B) -- (C) -- (D) -- (E);
\draw (F) -- (B);
\draw (D) -- (G);
\draw (D) -- (H);
\foreach \p in {A,B,C,D,E,F,G,H}{ \fill (\p) circle (0.055); }
% 虚线轮廓
\draw[dashed,rounded corners=26pt] (0.30,2.05) -- (1.00,2.95) -- (2.30,2.72) -- (3.05,1.55) -- (2.45,0.68) -- (0.75,1.10) -- cycle;
\node[font=\small,anchor=west] at (1.35,2.55) {非椭圆形结构};
\node[font=\small] at (1.7,-1.05) {(b) 链接效应};
\end{scope}

\end{tikzpicture}
\phantomsection\label{fig:12.5}\mvcap{图 12.5\quad 单连接聚类}
\end{mvfigure}

通过单连接法形成的聚类, 在对原始的距离(或相似性)作保序变换后保持不变. 具体说, 在\hyperref[tab:12.1]{表 12.1} 中, 任一组相互之间有单调关系的相似系数会产生相同的聚类.

\subsection{完全连接法}

完全连接聚类方式与单连接聚类方式在许多方面都相同, 但有一个重要区别: 在每个阶段, 聚类之间的距离(相似性)由两个聚类中相距最远的两个元素之间的距离确定. 这样, 完全
